% !TeX root = ../../Book1.tex
\chapter{Lebesgue 积分与收敛理论}
\label{b1:ch:05}
\BookChapterNumber{b1:ch:05}
% 本章资料及证明义务（2026-09-05，初稿，不计正式审核）：
% 5.1：Axler2020Measure §3A 74--80页；简单函数表示无关、无穷值约定、
% 递增简单逼近、非负可加、积分区域的可数可加，均不借尚未证明的一般MCT。
% 5.2：同书81--84、90--91、155--157页；有定义与可积区分、复线性、
% 等价类、简单逼近、绝对连续及有限测度局部化。
% 5.3：同书78、92--94页；Fatou用简单下界直接证明以避免纲要顺序造成循环；
% MCT、DCT、反向Fatou与Riemann相容性完整自证。
% 5.4：Axler DCT为工具来源；参数推论与具体计算独立推导。
% 5.5：Fremlin2016Measure vol.2 §246A/B/I/J官方全文，统一可积定义版本辨析；
% Bass2024Real 7.21--7.24及解答仅作背景对照，不列为正式正文引文。
% 中文伴读Zhou2016Shibianhanshulun只核验既有书目，不声称读到具体证明。
% 所有习题的来源或自拟方式另见题前注释。

\begin{chapterabstract}
上一章把可测函数分解、截断并逼近为简单函数；本章给这些逼近赋予可以取极限的数值。积分先在非负函数上定义，再通过正负部和实虚部推广到可积函数。单调收敛、Fatou 引理与控制收敛定理分别处理不同的极限过程。随后把这些结果用于含参变量积分，并先在有限测度空间用一致可积性说明：交换极限与积分需要控制的是积分质量，而不仅是每个点上的函数值；一般测度空间还须分开处理幅值尾部与空间尾部。除明确说明外，测度空间不预设有限或 \(\sigma\)-有限。
\end{chapterabstract}

\begin{learninggoals}
\begin{itemize}
\item\label{b1:item:0500-goal-1}能从简单函数的表示无关性出发完成积分构造，处理无穷值、零测集和可数分割。
\item\label{b1:item:0500-goal-2}能区分积分有定义与绝对可积，说明条件收敛级数或反常积分为何不能代替 Lebesgue 可积性，并解释 \(L^1\) 中为何把几乎处处相等的函数看成同一个元素。
\item\label{b1:item:0500-goal-3}能选择并证明所需的收敛定理，为参数极限或求导找到局部统一的可积控制。
\item\label{b1:item:0500-goal-4}能在有限测度空间用一致可积性与依测度收敛判断 \(L^1\) 收敛，并在一般测度空间区分幅值尾部控制与空间尾部控制。
\end{itemize}
\end{learninggoals}

本章的主要前置结果是\hyperref[b1:ch:02]{第二章}的测度连续性与\hyperref[b1:ch:04]{第四章}的可测运算、几乎处处约定和简单函数逼近。欧氏空间中的具体求积另用\hyperref[b1:ch:03]{第三章}的 Lebesgue 测度；一般积分理论本身不依赖欧氏几何。

\ProseParagraphBreak

先读前两节，确认积分是怎样构造出来的，再进入收敛定理。\secref{b1:sec:0501}只先证明递增简单函数的积分逼近；\secref{b1:sec:0503}将从 Fatou 引理推出一般单调收敛，避免在可加性的证明中预借尚未建立的结果。

\ProseParagraphBreak

\secref{b1:sec:0505}的一致可积判据及\cref{b1:ex:0508}的凸函数构造适合第二遍阅读。首次阅读可先掌握控制收敛及其参数应用，再完成\cref{b1:ex:0509}与\cref{b1:ex:0510}的基本辨析，最后回到第五节比较幅值尾部与空间尾部各自控制的误差。

\ProseParagraphBreak

Sheldon Axler 的《Measure, Integration \& Real Analysis》\cite[Chapter 3]{Axler2020Measure}适合伴读积分构造和极限定理，其中以可测分割起步的写法可以与本书的简单函数路线互相对照。Fremlin 的《Measure Theory》\cite[Volume 2, \S246]{Fremlin2016Measure}适合进一步核对一致可积性在一般测度空间中的定义与收敛判据；阅读时应注意其定义已经包含有限测度集合之外的统一控制。需要中文伴读时，可使用周民强的《实变函数论》\cite{Zhou2016Shibianhanshulun}对照 Lebesgue 积分与函数列收敛的传统讲法。

% !TeX root = ../../Book1.tex
\section{非负函数的积分}
\label{b1:sec:0501}
% 来源与证明义务：Axler2020Measure，§3A，印刷74--80页，已读全文。
% 改用简单函数起步；先证表示无关性、递增简单逼近和空间可数可加，
% 不前向调用5.3的一般单调收敛定理。共同分割与无穷值分支均展开。

测度已经规定了示性函数应有的积分：集合 \(A\) 的示性函数应当积成 \(\mu(A)\)。有限次加权由此给出简单函数的积分；对于一般非负函数，再取所有从下方逼近它的简单函数积分的上确界。构造中需要分别回答两个问题：同一个函数换一种表示，积分是否不变；从简单函数过渡到一般函数后，可加性是否仍然成立。

本节固定测度空间 \((X,\mathcal A,\mu)\)，不要求有限、\(\sigma\)-有限或完备。只在非负积分及其截断记号中约定 \(0\cdot\infty=\infty\cdot0=0\)。例如，零函数在无限测度空间上的积分仍为零；这项约定并不赋予一般扩展实值运算中的 \(\infty-\infty\) 任何意义。

\subsection{简单函数积分}
\label{b1:subsec:050101}

\begin{definition}\label{b1:def:simple-integral}
设非负简单函数具有规范表示
\[
s=\sum_{j=1}^{r}a_j\mathbf1_{A_j},
\quad a_j\in[0,\infty),
\]
其中 \((A_j)_{j=1}^{r}\) 是 \(X\) 的可测分割。定义它的\term[jiandan-hanshu-jifen]{简单函数积分}[integral of a simple function]为
\[
I_\mu(s)=\sum_{j=1}^{r}a_j\mu(A_j)\in[0,\infty].
\]
\end{definition}

规范表示的唯一性已见\cref{b1:prop:simple-canonical-form}，但计算时使用的表示未必规范。积分必须容许分块细分、合并以及集合重叠，不能把它们当作不同的函数。

\begin{proposition}\label{b1:prop:simple-integral-well-defined}
若 \(s=\sum_{j=1}^{r}a_j\mathbf1_{A_j}\)，其中 \(a_j\geq0\)、\(A_j\in\mathcal A\)，则不论这些集合是否互不相交，都有
\[
I_\mu(s)=\sum_{j=1}^{r}a_j\mu(A_j).
\]
非负简单函数的积分具有单调性，并且对 \(c,d\geq0\) 满足
\[
I_\mu(cs+dt)=cI_\mu(s)+dI_\mu(t).
\]
\end{proposition}
\begin{proof}
把每个 \(A_j\) 及其补集取交，得到有限可测分割 \((C_\ell)_\ell\)，略去空块。在每个 \(C_\ell\) 上，函数值为
\[
b_\ell=\sum_{\{j\mid C_\ell\subseteq A_j\}}a_j.
\]
每个规范水平集都是若干个这样的块之并。测度的有限可加性给出
\[
\begin{aligned}
I_\mu(s)
&=\sum_\ell b_\ell\mu(C_\ell)\\
&=\sum_{j=1}^{r}a_j\sum_{\{\ell\mid C_\ell\subseteq A_j\}}\mu(C_\ell)
=\sum_{j=1}^{r}a_j\mu(A_j).
\end{aligned}
\]
各项非负，有限次交换求和不涉及无穷量相减；零系数按照本节约定处理。

对两个简单函数取共同的有限可测分割，在每一块上分别记它们的值为 \(u_\ell,v_\ell\)。若 \(s\leq t\)，则 \(u_\ell\leq v_\ell\)，逐块乘以测度并求和便得单调性。函数 \(cs+dt\) 在该块上的值为 \(cu_\ell+dv_\ell\)，同一求和公式给出所述线性关系。若 \(c=0\) 或 \(d=0\)，对应零函数的积分为零，结论仍成立。
\end{proof}

例如，在实直线上取 \(s=2\mathbf1_{[0,2]}+3\mathbf1_{[1,3]}\)。两项的支集重叠，而积分仍为 \(2\cdot2+3\cdot2=10\)。也可以在 \([0,1)\)、\([1,2]\)、\((2,3]\) 上分别使用函数值 \(2,5,3\)，得到同一结果。这里用的是有限可测分割，不要求分块为区间。

\subsection{非负可测函数积分}
\label{b1:subsec:050102}

\begin{definition}\label{b1:def:nonnegative-integral}
设 \(f:X\rightarrow[0,\infty]\) 可测。它的\term[lebesgue-jifen]{Lebesgue 积分}[Lebesgue integral]定义为
\[
\int_X f\,\mathrm d\mu
=\sup\{\,I_\mu(s)\mid 0\leq s\leq f,\ s\in\mathcal S_+(X,\mathcal A)\,\},
\]
其中 \(\mathcal S_+(X,\mathcal A)\) 表示全部非负可测简单函数。当定义域明确时省略下标 \(X\)。对 \(E\in\mathcal A\)，定义
\[
\int_E f\,\mathrm d\mu=\int_X f\mathbf1_E\,\mathrm d\mu,
\]
其中 \(f\mathbf1_E\) 在 \(E\) 外取零，即使 \(f\) 在那里等于无穷也如此。
\symidx{\(\int_E f\,\mathrm d\mu\)：函数关于测度的积分}
\end{definition}

候选集合至少包含零函数，因此上确界总有意义。若 \(f\) 本身简单，简单函数积分的单调性说明每个候选的积分不超过 \(I_\mu(f)\)，而 \(f\) 又是候选，所以新定义恰好等于旧定义。另一方面，\(f\leq g\) 时，\(f\) 的每个候选也是 \(g\) 的候选，故新积分同样保序。

下面的逼近引理把上确界变成一个数列极限。这里只处理递增的简单函数列；它将在建立可加性时使用，不预借后面的收敛定理。

\begin{lemma}\label{b1:lem:simple-integral-approximation}
若非负简单函数列满足 \(s_n\uparrow f\)，则
\[
I_\mu(s_n)\uparrow\int f\,\mathrm d\mu.
\]
因此这个极限与所选的递增简单逼近列无关。
\end{lemma}
\begin{proof}
由单调性，\(L=\lim_n I_\mu(s_n)\) 存在且不超过 \(\int f\,\mathrm d\mu\)。固定任意非负简单函数 \(s\leq f\)，写成 \(s=\sum_ja_j\mathbf1_{A_j}\)，并固定 \(0<c<1\)。令
\[
E_n=\{\,x\in X\mid s_n(x)\geq cs(x)\,\}.
\]
这些集合递增，且并为 \(X\)：当 \(s(x)=0\) 时每个不等式都成立；当 \(s(x)>0\) 时，\(s_n(x)\to f(x)\geq s(x)>cs(x)\)，不等式最终成立。于是
\[
I_\mu(s_n)\geq cI_\mu(s\mathbf1_{E_n})
=c\sum_j a_j\mu(A_j\cap E_n).
\]
测度的从下连续性给出 \(L\geq cI_\mu(s)\)。若 \(I_\mu(s)=\infty\)，这已经说明 \(L=\infty\)；否则令 \(c\uparrow1\)，得 \(L\geq I_\mu(s)\)。最后对所有 \(s\leq f\) 取上确界，便得反向不等式。
\end{proof}

\Cref{b1:thm:nonnegative-simple-approximation}保证每个非负可测函数都有这样的逼近列，故引理适用于全部积分对象。例如函数在正测度集合上等于 \(+\infty\) 时，可以在该集合上取任意高的常值简单函数，其积分必为无穷。无穷函数值若只出现在零测集上，则情况不同，下面会给出精确结论。

\subsection{单调性}
\label{b1:subsec:050103}

\begin{proposition}\label{b1:prop:nonnegative-integral-linearity}
对非负可测函数 \(f,g\) 以及有限常数 \(c,d\geq0\)，有
\[
\int(cf+dg)\,\mathrm d\mu
=c\int f\,\mathrm d\mu+d\int g\,\mathrm d\mu.
\]
此外，\(f\leq g\) 蕴含 \(\int f\,\mathrm d\mu\leq\int g\,\mathrm d\mu\)。
\end{proposition}
\begin{proof}
取非负简单函数 \(s_n\uparrow f\)、\(t_n\uparrow g\)。当 \(c,d>0\) 时，\(cs_n+dt_n\uparrow cf+dg\)，由简单函数的线性与\cref{b1:lem:simple-integral-approximation}，
\[
\int(cf+dg)\,\mathrm d\mu
=\lim_n\bigl(cI_\mu(s_n)+dI_\mu(t_n)\bigr)
=c\int f\,\mathrm d\mu+d\int g\,\mathrm d\mu.
\]
非负数列的和允许一项或两项趋于无穷，因而最后一步仍成立。零系数的情形直接略去对应项。单调性已由定义中的候选集合包含关系得到。
\end{proof}

\begin{proposition}\label{b1:prop:integral-markov}
设 \(f\geq0\) 可测。对 \(a>0\)，有
\[
a\mu(\{f\geq a\})\leq\int f\,\mathrm d\mu.
\]
并且，\(\int f\,\mathrm d\mu=0\) 当且仅当 \(f=0\) 几乎处处；若积分有限，则 \(f\) 几乎处处有限。
\end{proposition}
\begin{proof}
由 \(a\mathbf1_{\{f\geq a\}}\leq f\) 及单调性，得到第一个估计。若积分为零，则对每个正整数 \(k\)，集合 \(\{f\geq1/k\}\) 为零测集。它们的并为 \(\{f>0\}\)，故 \(f=0\) 几乎处处。

反过来，设 \(f\) 在可测零集 \(N\) 外为零。每个非负简单函数 \(s\leq f\) 的正水平集都包含于 \(N\)，从而各项积分均为零；取上确界即得结论。若 \(\int f\,\mathrm d\mu=M<\infty\)，则
\[
\mu(\{f=\infty\})\leq\mu(\{f\geq k\})\leq M/k
\]
对每个 \(k\) 成立，令 \(k\to\infty\) 得最后一个断言。
\end{proof}

\begin{corollary}\label{b1:cor:integral-ae-invariance}
若非负可测函数 \(f,g\) 几乎处处相等，则它们的积分相等。积分单调性中的 \(f\leq g\) 也可以只要求几乎处处成立。
\end{corollary}
\begin{proof}
取包含例外点的可测零集 \(N\)。分解 \(f=f\mathbf1_{X\setminus N}+f\mathbf1_N\)，再用非负可加性及零积分判别，得到
\(\int f=\int f\mathbf1_{X\setminus N}\)。对 \(g\) 同样处理，在补集上使用相等或大小关系即可。
\end{proof}

这里始终要求参与积分的函数可测。在非完备空间中，不能因为任意修改只发生在一个零测集上，就不再检查修改后的函数是否可测。对于 Lebesgue 测度，完备性允许这样的修改，例如 \(\mathbf1_{\mathbb Q\cap[0,1]}\) 的积分为零，而 \(\mathbf1_{[0,1]\setminus\mathbb Q}\) 的积分为一。

\subsection{可数可加性}
\label{b1:subsec:050104}

固定函数，把积分看成积分区域的函数，会重新得到一个测度。这一步说明积分保留了原先测度的可数结构。

\begin{theorem}\label{b1:thm:integral-induced-measure}
设 \(f:X\rightarrow[0,\infty]\) 可测。公式
\[
\nu_f(E)=\int_E f\,\mathrm d\mu\quad(E\in\mathcal A)
\]
定义了 \((X,\mathcal A)\) 上的一个测度。因此，对两两不交的可测集 \((E_j)_{j\geq1}\)，
\[
\int_{\bigcup_jE_j}f\,\mathrm d\mu
=\sum_{j=1}^{\infty}\int_{E_j}f\,\mathrm d\mu.
\]
\end{theorem}
\begin{proof}
先对简单函数 \(s=\sum_\ell a_\ell\mathbf1_{A_\ell}\) 考虑
\[
\nu_s(E)=\sum_\ell a_\ell\mu(A_\ell\cap E).
\]
它是有限个测度的非负线性组合，空集上为零，可数可加性由 \(\mu\) 的可数可加性逐项得到。

对一般 \(f\) 取 \(s_n\uparrow f\)。在每个 \(E\) 上，\(s_n\mathbf1_E\uparrow f\mathbf1_E\)，所以\cref{b1:lem:simple-integral-approximation}给出
\(\nu_f(E)=\lim_n\nu_{s_n}(E)\)，且这个极限递增。令 \(E=\bigcup_jE_j\)。一方面，
\[
\nu_{s_n}(E)=\sum_j\nu_{s_n}(E_j)\leq\sum_j\nu_f(E_j),
\]
令 \(n\to\infty\) 得到一个方向。另一方面，对每个有限的 \(r\)，
\[
\nu_{s_n}(E)\geq\sum_{j=1}^{r}\nu_{s_n}(E_j).
\]
先令 \(n\to\infty\)，再令 \(r\to\infty\)，便得另一个方向。全过程只用了有限个递增非负数列的极限与非负级数的定义。
\end{proof}

取 \(f=1\) 就恢复原测度。若 \(f\) 非负且积分为一，\(\nu_f\) 是概率测度。其反问题是：什么测度能够这样由函数生成？\hyperref[b1:ch:06]{下一章}将回答这个问题。

在正整数集上取计数测度 \(\#\)。单点 \(\{k\}\) 上的积分等于函数值，故本定理给出
\[
\int_{\mathbb N}f\,\mathrm d\#=\sum_{k=1}^{\infty}f(k)\quad(f\geq0).
\]
积分因而把非负级数与连续区域上的求积纳入同一规则。稍后对函数列求和的定理，将进一步允许在适当条件下把求和号移出积分号。

% !TeX root = ../../Book1.tex
\section{一般可积函数}
\label{b1:sec:0502}

% TODO: 撰写本节正文。

\subsection{正部与负部}
\label{b1:subsec:050201}

待填充。

\subsection{可积性}
\label{b1:subsec:050202}

待填充。

\subsection{\texorpdfstring{\(L^1\)}{L1} 的初步定义}
\label{b1:subsec:050203}

待填充。

\subsection{积分的绝对连续性}
\label{b1:subsec:050204}

待填充。

% !TeX root = ../../Book1.tex
\section{基本收敛定理}
\label{b1:sec:0503}
% 主来源：Axler2020Measure §3A 3.11、§3B 3.31及3.34，印刷78、92--94页。
% Fatou按简单下界和递增集合直接证明，以遵守纲要顺序而不循环引用MCT。
% Riemann相容性只证本章计算所需方向；连续性几乎处处的充要判别不在此展开。

点态极限只控制每个固定点上的函数值，不控制积分在整个空间中的分布。以
\[
f_n=n\mathbf1_{(0,1/n)}
\]
为例，它在 \((0,1)\) 上处处趋于零，积分却恒为一。高度增长抵消了支集缩小。收敛定理的假设应当排除这种损失，或至少规定极限积分与原积分之间的不等式。

\subsection{Fatou 引理}
\label{b1:subsec:050301}

对于非负函数，即使没有共同上界，仍有一个方向的不等式。\term*[fatou-yinli]{Fatou 引理}[Fatou's lemma]允许极限丢失质量，却不允许非负质量在逐点极限中凭空增加。

\begin{lemma}\label{b1:thm:fatou}
设 \((f_n)\) 是测度空间上的非负可测函数列，则
\[
\int\liminf_{n\to\infty}f_n\,\mathrm d\mu
\leq\liminf_{n\to\infty}\int f_n\,\mathrm d\mu.
\]
\end{lemma}
\begin{proof}
记 \(h=\liminf_n f_n\)，它由可数极限运算的可测性仍为非负可测函数。固定非负简单函数 \(s\leq h\) 以及 \(0<c<1\)。对每个 \(k\) 令
\[
E_k=\bigcap_{n\geq k}\{\,x\in X\mid f_n(x)\geq cs(x)\,\}.
\]
集合 \(E_k\) 可测且递增。若 \(s(x)=0\)，则 \(x\) 属于全部 \(E_k\)；若 \(s(x)>0\)，由 \(\liminf_n f_n(x)\geq s(x)>cs(x)\)，知 \(x\) 最终也属于 \(E_k\)。因此 \(E_k\uparrow X\)。

固定 \(k\)，当 \(n\geq k\) 时有 \(f_n\geq cs\mathbf1_{E_k}\)，于是
\[
\liminf_n\int f_n\,\mathrm d\mu\geq c\int_{E_k}s\,\mathrm d\mu.
\]
右端是简单函数积分；对其有限个水平集使用测度从下连续性，令 \(k\to\infty\)，得到右端趋于 \(cI_\mu(s)\)。若这个积分无穷，所需下界已得；否则再令 \(c\uparrow1\)。最后对全部 \(s\leq h\) 取上确界，便得到结论。
\end{proof}

章首的尖峰列使不等号严格。非负性也不能直接删去：取 \(-f_n\)，极限积分仍为零，而各项积分为负一。若存在共同可积控制，就可以通过非负函数重新利用 Fatou 引理；控制收敛的证明正采用这种办法。

\subsection{单调收敛定理}
\label{b1:subsec:050302}

非负递增列没有先出现、后消失的尖峰。它把积分的下逼近结构保留下来，因而 Fatou 不等式能够加强为等式。

\begin{theorem}\label{b1:thm:monotone-convergence}
设非负可测函数满足 \(f_n\leq f_{n+1}\) 几乎处处，并且 \(f_n\to f\) 几乎处处，其中 \(f:X\rightarrow[0,\infty]\) 可测。则
\[
\int f_n\,\mathrm d\mu\uparrow\int f\,\mathrm d\mu.
\]
这称为\term[dandiao-shoulian-dingli]{单调收敛定理}[monotone convergence theorem]。
\end{theorem}
\begin{proof}
把单调性与收敛性的可数个例外零集取并，并将全部函数在所得可测零集上改为零。分块定义保证修改后的函数可测，积分由\cref{b1:cor:integral-ae-invariance}保持不变。因此只需证明逐点版本。

记 \(L=\lim_n\int f_n\,\mathrm d\mu\)。它由单调性存在，而 \(f_n\leq f\) 给出 \(L\leq\int f\,\mathrm d\mu\)。Fatou 引理给出反向不等式
\(\int f\,\mathrm d\mu\leq\liminf_n\int f_n\,\mathrm d\mu=L\)，两者合并即得等式。
\end{proof}

\begin{corollary}\label{b1:cor:nonnegative-series-integral}
对任意非负可测函数列 \((u_j)\)，有
\[
\int\sum_{j=1}^{\infty}u_j\,\mathrm d\mu
=\sum_{j=1}^{\infty}\int u_j\,\mathrm d\mu.
\]
\end{corollary}
\begin{proof}
部分和 \(S_n=\sum_{j=1}^{n}u_j\) 递增至非负级数的和。先对有限和使用\cref{b1:prop:nonnegative-integral-linearity}，再应用单调收敛定理，便得到等式。两边允许同时为无穷。
\end{proof}

这条结果不要求级数预先可积。例如在 \((0,1)\) 上，几何级数给出
\[
\frac1{1-x}=\sum_{j=0}^{\infty}x^j.
\]
一旦核定初等积分与本节积分的一致性，便可逐项积分，得到发散的调和级数。此时结论是原函数不可积，而不是“逐项积分失败”。

\begin{corollary}\label{b1:cor:decreasing-integral-convergence}
若 \(0\leq f_n\downarrow f\) 几乎处处，且 \(\int f_1\,\mathrm d\mu<\infty\)，则
\[
\int f_n\,\mathrm d\mu\downarrow\int f\,\mathrm d\mu.
\]
\end{corollary}
\begin{proof}
略去共同可测零集后，还可令各函数处处有限。函数 \(f_1-f_n\) 非负递增至 \(f_1-f\)，由单调收敛定理及有限积分的线性，得到
\[
\int f_1\,\mathrm d\mu-\int f_n\,\mathrm d\mu
\longrightarrow
\int f_1\,\mathrm d\mu-\int f\,\mathrm d\mu.
\]
减去共同的有限量即可。
\end{proof}

有限起始积分不可省。在实直线上，\(\mathbf1_{[n,\infty)}\downarrow0\)，但每个积分都为无穷。这与测度从上连续性需要一个有限测度起点完全对应。

\subsection{控制收敛定理}
\label{b1:subsec:050303}

积分运算中最常见的情形既不非负，也不单调。一个与序号无关的可积函数若能控制全部函数的绝对值，就足以交换极限与积分。Sheldon Axler 的《Measure, Integration \& Real Analysis》\cite{Axler2020Measure}在这一主题中还展示了 Egoroff 定理与积分绝对连续性配合的证明；下面采用 Fatou 引理的路线，直接得到绝对积分误差趋零。

\begin{theorem}\label{b1:thm:dominated-convergence}
设 \(f_n,f\) 是实值或复值可测函数，\(f_n\to f\) 几乎处处。若存在非负可积函数 \(g\)，使每个 \(n\) 都满足 \(|f_n|\leq g\) 几乎处处，则 \(f_n,f\) 均可积，而且
\[
\lim_{n\to\infty}\int|f_n-f|\,\mathrm d\mu=0,
\quad
\lim_{n\to\infty}\int f_n\,\mathrm d\mu=\int f\,\mathrm d\mu.
\]
这称为\term[kongzhi-shoulian-dingli]{控制收敛定理}[dominated convergence theorem]。
\end{theorem}
\begin{proof}
在一个共同可测零集外，收敛与所有控制不等式同时成立，且 \(g\) 有限。取极限知 \(|f|\leq g\)，由积分单调性得到可积性。令
\[
h_n=2g-|f_n-f|
\]
并在例外零集上取零。这是非负可测函数，且几乎处处趋于 \(2g\)。对它应用 Fatou 引理，
\[
2\int g\,\mathrm d\mu
\leq\liminf_n\left(2\int g\,\mathrm d\mu-\int|f_n-f|\,\mathrm d\mu\right).
\]
共同量 \(2\int g\,\mathrm d\mu\) 有限，所以
\(\limsup_n\int|f_n-f|\,\mathrm d\mu\leq0\)。被积函数非负，第一项极限即为零；再用\cref{b1:prop:integral-triangle}，
\[
\left|\int f_n\,\mathrm d\mu-\int f\,\mathrm d\mu\right|
\leq\int|f_n-f|\,\mathrm d\mu,
\]
得到第二项极限。
\end{proof}

有限测度空间上，若 \(|f_n|\leq M\) 几乎处处，可以取 \(g=M\mathbf1_X\)，这就是有界收敛的情形。无限测度空间上常数函数未必可积，所以仅有一致有界不够。\(\mathbf1_{[n,n+1]}\) 在实直线上处处趋零，而积分恒为一，正是一个反例。

\begin{proposition}\label{b1:prop:riemann-lebesgue-compatibility}
若有界函数 \(f:[a,b]\rightarrow\mathbb R\) Riemann 可积，则它 Lebesgue 可测且可积，并且两种积分相等。
\end{proposition}
\begin{proof}
取逐次加细的有限分割 \(P_n\)，使上、下 Darboux 和之差小于 \(1/n\)。记相应的下、上阶梯函数为 \(l_n,u_n\)，在分割端点处的值任意取为统一有界的值。所有分割端点组成一个可数零测集 \(D\)。在其补集上，
\[
-M\leq l_n\leq l_{n+1}\leq f\leq u_{n+1}\leq u_n\leq M,
\]
其中 \(M\) 是 \(f\) 的绝对值上界。可测极限 \(l=\lim l_n\)、\(u=\lim u_n\) 在 \(D\) 上另取零；由于 \(D\) 可测，分块定义仍可测。控制收敛给出
\[
\int_{[a,b]}(u-l)\,\mathrm dm
=\lim_n\int_{[a,b]}(u_n-l_n)\,\mathrm dm=0.
\]
于是 \(l=u\) 几乎处处，而 \(l\leq f\leq u\) 在 \(D\) 外成立。Lebesgue 测度的完备性说明 \(f\) 可测；有界性与区间有限测度又给出可积性。最后 \(l_n\to f\) 几乎处处，再用控制收敛，得到 \(f\) 的 Lebesgue 积分等于下 Darboux 和的极限，也就是它的 Riemann 积分。
\end{proof}

因此连续函数在紧区间上的初等求积公式可以直接使用。对非负函数还可由单调收敛令区间递增，例如
\[
\int_0^1 x^{-p}\,\mathrm dx=
\begin{cases}
(1-p)^{-1},&p<1,\\
+\infty,&p\geq1.
\end{cases}
\]
在 \([\varepsilon,1]\) 上先算初等积分，再令 \(\varepsilon\downarrow0\)，便得到该式。这里近零端的可积性由奇性阶数决定，与远处尾部的可积性是两种不同检查。

\begin{corollary}\label{b1:cor:absolute-series-integral}
若实值或复值可积函数列满足 \(\sum_j\int|u_j|\,\mathrm d\mu<\infty\)，则级数 \(\sum_ju_j\) 几乎处处绝对收敛，其和可取可积代表 \(S\)，并且
\[
\int S\,\mathrm d\mu=\sum_{j=1}^{\infty}\int u_j\,\mathrm d\mu.
\]
\end{corollary}
\begin{proof}
由非负级数的积分公式，\(G=\sum_j|u_j|\) 的积分有限，故 \(G<\infty\) 几乎处处。在这个可测集合上定义级数和，在补集上令和为零，得到可测函数 \(S\)。部分和几乎处处受 \(G\) 控制并趋于 \(S\)，应用控制收敛即可。数值级数也绝对收敛，因为 \(\sum_j|\int u_j|\leq\sum_j\int|u_j|<\infty\)。
\end{proof}

\subsection{反向 Fatou 型结论}
\label{b1:subsec:050304}

如果函数列有共同可积上界，可以反转 Fatou 不等式的方向；此时起作用的仍是非负差函数。

\begin{proposition}\label{b1:prop:reverse-fatou}
设 \(f_n\) 为实值可积函数，存在实值可积函数 \(g\)，使 \(f_n\leq g\) 几乎处处。则 \(h=\limsup_n f_n\) 的积分有定义，允许为 \(-\infty\)，并且
\[
\limsup_n\int f_n\,\mathrm d\mu\leq\int h\,\mathrm d\mu.
\]
\end{proposition}
\begin{proof}
在共同零测集外有 \(h\leq g\)，故 \(h^+\leq g^+\)，所以 \(\int h^+<\infty\)，积分不会成为 \(\infty-\infty\)。对非负函数 \(g-f_n\) 应用 Fatou 引理，得到
\[
\int(g-h)\,\mathrm d\mu
\leq\int g\,\mathrm d\mu-\limsup_n\int f_n\,\mathrm d\mu.
\]
其中 \(\liminf_n(g-f_n)=g-h\)，在 \(h=-\infty\) 处差取 \(+\infty\)。为说明左端的展开合法，可用非负恒等式
\[
(g-h)+g^-+h^+=g^++h^-
\]
积分；除 \(g-h\) 与 \(h^-\) 外，其余项积分均有限。因而
\(\int(g-h)=\int g-\int h\)，包括两边对应无穷的情形。代入上式并移去有限量 \(\int g\)，就得到所需不等式。
\end{proof}

对非负尖峰列 \(n\mathbf1_{(0,1/n)}\)，反向不等式会要求 \(1\leq0\)，所以它不可能有共同可积上界。使用本节结果时，可以先检查手中的控制是哪一种：单调性给精确的非负极限，单侧可积控制给一个方向的不等式，双侧绝对值控制则给出 \(L^1\) 收敛。

% !TeX root = ../../Book1.tex
\section{参数积分：极限与微分}
\label{b1:sec:0504}

% 资料：已核对 Axler2020Measure，第 3B 节，3.28--3.31，印刷页 90--93。
% 该来源用于控制收敛及可积函数的局部化；以下参数推论由本章结果自证，
% 不将参数求导定理或计算实例归作 Axler 原书的直接陈述。
% 证明义务：参数邻域内统一控制；变动区间的端点零集；公共零集外的
% 参数可微性、导数可测性与差商控制；先截断再处理无穷远的一致尾部。

控制收敛定理也能处理连续变化的参数。设 \((X,\mathcal A,\mu)\) 是测度空间，实参数 \(t\) 在区间 \(I\) 中变化。若每个截面 \(x\mapsto f(t,x)\) 可积，便可定义\term[canshu-jifen]{参数积分}[parameter integral]，记为
\[
F(t)=\int_X f(t,x)\,\mathrm d\mu(x).
\]
逐点检查 \(f\) 对参数的变化以后，还要找到对附近参数同时有效的可积控制。控制函数在这里的作用，可对照 Sheldon Axler 的《Measure, Integration \& Real Analysis》\cite{Axler2020Measure}第 3B 节关于控制收敛的论证。以下允许被积函数取实值或复值；记号 \(dx\) 表示对一维 Lebesgue 测度积分。

\subsection{参数积分的连续性}
\label{b1:subsec:050401}

\begin{proposition}
\label{b1:prop:parameter-integral-continuity}
设 \(t_0\in I\)，每个 \(f(t,\cdot)\) 可测。若存在 \(t_0\) 在 \(I\) 中的邻域 \(U\)、非负可积函数 \(g\) 和可测零集 \(N\)，使对所有 \(x\notin N\) 都有
\[
\lim_{t\to t_0}f(t,x)=f(t_0,x),
\quad
|f(t,x)|\leq g(x)\quad(t\in U),
\]
则 \(F\) 在 \(U\) 上有定义，并且在 \(t_0\) 处连续。
\end{proposition}
\begin{proof}
可积控制先保证各个积分有限。任取 \(U\) 中趋于 \(t_0\) 的序列 \(t_n\)，由\cref{b1:thm:dominated-convergence}，
\[
|F(t_n)-F(t_0)|
\leq\int_X|f(t_n,x)-f(t_0,x)|\,\mathrm d\mu(x)
\longrightarrow0.
\]
实数区间上的连续性可由序列刻画，故结论成立。
\end{proof}

证明还给出了截面在 \(L^1(\mu)\) 中的连续性。条件只放在 \(t_0\) 的一个邻域上，因此研究另一个参数点时可以换用新的控制函数，不必用同一个上界控制整个参数区间。每个截面各自可积却不够：这只说明积分存在，并未限制附近截面的大小如何变化。

例如，若 \(h\in L^1(\mathbb R)\)，则
\[
t\longmapsto\int_{\mathbb R}e^{itx}h(x)\,\mathrm dx
\]
连续，因为被积函数的绝对值恒为 \(|h(x)|\)。这里没有要求 \(xh(x)\) 可积；若要再对参数求导，就需要检查差商能否受到可积函数控制。

\subsection{积分号下取极限}
\label{b1:subsec:050402}

参数的极限点不必属于原来的参数范围。例如研究 \(t\downarrow0\) 时，可把逐点极限作为 \(t=0\) 的截面，再应用\cref{b1:prop:parameter-integral-continuity}。积分区间也可以随参数改变，只要先把它写进示性函数。

\begin{proposition}
\label{b1:prop:moving-interval-integral-limit}
设 \(a_n,b_n\in[A,B]\)，\(a_n\leq b_n\)，且 \(a_n\to a\)、\(b_n\to b\)。设 \(f_n,f\) 在 \([A,B]\) 上可测，\(f_n\to f\) 几乎处处，并有非负 \(g\in L^1([A,B])\) 使 \(|f_n|\leq g\) 几乎处处。则
\[
\lim_{n\to\infty}\int_{a_n}^{b_n}f_n(x)\,\mathrm dx
=\int_a^b f(x)\,\mathrm dx.
\]
\end{proposition}
\begin{proof}
除端点 \(a,b\) 外，示性函数 \(\mathbf1_{[a_n,b_n]}\) 逐点收敛于 \(\mathbf1_{[a,b]}\)。这两个端点组成零集；再合并函数列的可数个例外零集，即得
\[
f_n\mathbf1_{[a_n,b_n]}
\longrightarrow f\mathbf1_{[a,b]}
\quad\text{几乎处处}.
\]
左侧仍由 \(g\) 控制，在固定区间 \([A,B]\) 上应用\cref{b1:thm:dominated-convergence}即可。
\end{proof}

端点的零测性在这里承担了具体作用。若把 Lebesgue 测度换成集中在 \(0\) 的单位质量，则 \([1/n,1]\) 的积分恒为零，而极限区间 \([0,1]\) 的积分为一。一般测度下，要么另行控制端点的质量，要么直接检查变动集合的示性函数是否几乎处处收敛。

取 \(0<|t|<1/2\)。在 \([0,2]\) 上有 \(|\sin(tx)/t|\leq x\)，且该商趋于 \(x\)，所以
\[
\lim_{t\to0}\int_0^{1+t}\dfrac{\sin(tx)}{t}\,\mathrm dx
=\int_0^1x\,\mathrm dx=\dfrac12.
\]
这里只需对任意趋于零的参数序列应用\cref{b1:prop:moving-interval-integral-limit}。末项及下文连续函数的有限区间积分，可以依据\cref{b1:prop:riemann-lebesgue-compatibility}用初等定积分计算。

\subsection{积分号下求导}
\label{b1:subsec:050403}

求导要求控制的对象是差商。导数的一致上界可以通过一元中值定理给出这种控制，但例外零集必须与参数无关。

\begin{theorem}[积分号下求导]
\label{b1:thm:differentiation-under-integral}
设 \(J\) 是开区间，每个 \(f(t,\cdot)\) 可测，且某个 \(t_*\in J\) 满足 \(f(t_*,\cdot)\in L^1(\mu)\)。设存在可测零集 \(N\) 和非负 \(g\in L^1(\mu)\)，使对每个 \(x\notin N\)，函数 \(t\mapsto f(t,x)\) 在整个 \(J\) 上可微，并满足
\[
|\partial_t f(t,x)|\leq g(x)\quad(t\in J).
\]
在 \(N\) 上把导数定义为零。则每个 \(f(t,\cdot)\) 可积，\(F\) 在 \(J\) 上可微，且
\[
F'(t)=\int_X\partial_t f(t,x)\,\mathrm d\mu(x).
\]
\end{theorem}
\begin{proof}
先考虑实值函数。对 \(x\notin N\)，中值定理给出
\[
|f(t,x)-f(s,x)|\leq |t-s|g(x)
\quad(s,t\in J).
\]
复值情形分别对实部、虚部使用中值定理，再用三角不等式，得到右侧改为 \(2|t-s|g(x)\) 的估计。因此两种情形都满足
\[
|f(t,x)|\leq |f(t_*,x)|+2|t-t_*|g(x),
\]
从而每个截面可积。

固定 \(t\in J\)，取非零实数 \(h_n\to0\) 且 \(t+h_n\in J\)。在 \(X\setminus N\) 上，相应差商趋于 \(\partial_t f(t,x)\)，绝对值不超过 \(2g(x)\)。把差商在 \(N\) 上置零，就得到可测函数列，故其极限即上述导数也是可测的。由\cref{b1:thm:dominated-convergence}与积分线性性，
\[
\dfrac{F(t+h_n)-F(t)}{h_n}
=\int_X\dfrac{f(t+h_n,x)-f(t,x)}{h_n}\,\mathrm d\mu(x)
\longrightarrow\int_X\partial_t f(t,x)\,\mathrm d\mu(x).
\]
该结论对任意这样的序列成立，因而给出 \(F'(t)\)。
\end{proof}

基准截面的可积性也不能删去。在 \((0,\infty)\) 上取 \(f(t,x)=1+te^{-x}\)，其参数导数由可积函数 \(e^{-x}\) 控制，但每个截面都不可积。证明中先用基准截面控制函数本身，再用导数控制差商，两处可积性分别保证 \(F\) 有定义和极限可以交换。

若只要求对每个固定参数几乎处处可微，结论可能失败。在 \(I=X=(0,1)\) 上取 \(f(t,x)=\mathbf1_{(0,t)}(x)\)。对固定 \(t\)，除 \(x=t\) 外的参数导数都是零，但 \(F(t)=t\) 的导数是 \(1\)。不同参数的例外点合起来覆盖了整个区间，不能从这里删去一个公共零集。

\begin{corollary}
\label{b1:cor:higher-parameter-derivatives}
设 \(I\) 是开区间。在一个公共可测零集外，\(f(t,x)\) 关于 \(t\in I\) 属于 \(C^k\)，每个 \(f(t,\cdot)\) 可测；各正阶导数在该零集上取零。若每个参数点都有开区间邻域 \(J\subseteq I\)，使各阶 \(0\leq j\leq k\) 的导数在 \(J\) 上分别受到某个非负可积函数 \(g_j\) 控制，则 \(F\in C^k(I)\)，且
\[
F^{(j)}(t)=\int_X\partial_t^j f(t,x)\,\mathrm d\mu(x)
\quad(0\leq j\leq k).
\]
\end{corollary}
\begin{proof}
逐阶取差商，得到各阶导数截面的可测性。固定上述邻域 \(J\)，第零阶控制保证可积；依次对 \(f,\partial_t f,\ldots,\partial_t^{k-1}f\) 应用\cref{b1:thm:differentiation-under-integral}，得到各阶公式。每阶被积函数关于参数连续且有可积控制，\cref{b1:prop:parameter-integral-continuity}又保证各个 \(F^{(j)}\) 连续。
\end{proof}

回到本节开始的指数积分。若除 \(h\in L^1(\mathbb R)\) 外，还有 \(|x|^kh(x)\in L^1(\mathbb R)\)，则每个 \(j\leq k\) 的参数导数都由 \((1+|x|^k)|h(x)|\) 控制。这是因为 \(|x|^j\leq1+|x|^k\)。于是该参数积分属于 \(C^k\)，并且
\[
\dfrac{d^j}{dt^j}\int_{\mathbb R}e^{itx}h(x)\,\mathrm dx
=\int_{\mathbb R}(ix)^j e^{itx}h(x)\,\mathrm dx.
\]
这些加权可积条件把空间变量的增长与参数方向的可微性联系起来。

\subsection{无穷区间积分中的参数依赖}
\label{b1:subsec:050404}

在无穷区间上，有限区间内的控制还不足以决定整个积分的极限。截断以后的尾部需要对附近参数同时变小。

\begin{proposition}
\label{b1:prop:parameter-uniform-tails}
设 \(f(t,\cdot)\in L^1(0,\infty)\)，\(U\) 是 \(t_0\) 的参数邻域。若对每个 \(R>0\) 都有
\[
\lim_{t\to t_0}\int_0^R|f(t,x)-f(t_0,x)|\,\mathrm dx=0,
\]
并且
\[
\lim_{R\to\infty}\sup_{t\in U}\int_R^\infty|f(t,x)|\,\mathrm dx=0,
\]
则 \(f(t,\cdot)\to f(t_0,\cdot)\) 在 \(L^1(0,\infty)\) 中成立，参数积分也趋于 \(F(t_0)\)。
\end{proposition}
\begin{proof}
把差的绝对值积分拆成 \((0,R)\) 和 \((R,\infty)\) 两段，得到
\[
\begin{split}
\int_0^\infty|f(t,x)-f(t_0,x)|\,\mathrm dx
&\leq\int_0^R|f(t,x)-f(t_0,x)|\,\mathrm dx\\
&\quad+2\sup_{s\in U}\int_R^\infty|f(s,x)|\,\mathrm dx.
\end{split}
\]
给定 \(\varepsilon>0\)，先选 \(R\) 使第二项小于 \(\varepsilon/2\)，再令 \(t\) 足够接近 \(t_0\) 使第一项小于 \(\varepsilon/2\)。积分的收敛由\cref{b1:prop:integral-triangle}得到。
\end{proof}

第一项常可在紧矩形上处理。例如，若 \(f\) 在 \([t_0-\delta,t_0+\delta]\times[0,R]\) 上连续，则它在这个紧集上有界，常数上界在有限区间上可积，控制收敛便给出命题中的局部条件。无穷远的部分仍须另作估计。选择截断位置时，应先使尾部对整个参数邻域都小，再固定这个位置考察参数极限；逐个参数选择不同的截断位置，不能替代所需的一致估计。

若附近所有截面都由同一个 \(g\in L^1(0,\infty)\) 控制，则\cref{b1:thm:monotone-convergence}应用于 \(g\mathbf1_{(0,R)}\) 就给出一致尾部条件。反之，区间族 \(f(t,x)=\mathbf1_{[t,t+1]}(x)\) 在 \(t\to\infty\) 时逐点趋于零，在任意固定有限区间内最终恒为零，而整个积分始终为 \(1\)。它的质量移向无穷远，尾部不能一致舍去。

\begin{proposition}
\label{b1:prop:laplace-exponential-moments}
对 \(a>0\) 和整数 \(m\geq0\)，有
\[
\int_0^\infty x^m e^{-ax}\,\mathrm dx=\dfrac{m!}{a^{m+1}}.
\]
函数 \(L(a)=\int_0^\infty e^{-ax}\,\mathrm dx\) 在 \((0,\infty)\) 上光滑，其第 \(m\) 阶导数可在积分号内计算。
\end{proposition}
\begin{proof}
先在有限区间计算，再用单调收敛，得到
\[
L(a)=\lim_{R\to\infty}\dfrac{1-e^{-aR}}{a}=\dfrac1a.
\]
固定 \(a_0>0\)。当 \(a_0/2<a<3a_0/2\) 时，每阶参数导数满足
\[
|\partial_a^m e^{-ax}|=x^m e^{-ax}
\leq C_{m,a_0}e^{-a_0x/4},
\]
其中可取 \(C_{m,a_0}=m!(4/a_0)^m\)：把初等不等式 \(e^y\geq y^m/m!\) 用于 \(y=a_0x/4\)，就得到所需上界。右侧由已经算出的零阶积分可知可积。\cref{b1:cor:higher-parameter-derivatives}于是适用于任意有限阶数，并给出
\[
(-1)^m\int_0^\infty x^m e^{-ax}\,\mathrm dx
=L^{(m)}(a)=\dfrac{(-1)^m m!}{a^{m+1}}.
\]
\end{proof}

这里的控制只需在每个 \(a_0>0\) 的邻域内成立；当 \(a\downarrow0\) 时，\(L(a)\to\infty\)，不能把同一结论延伸到零参数。第\ref{b1:sec:0505}节将进一步讨论：一个函数族没有共同的可积上界时，还能用哪些条件控制积分极限。

% !TeX root = ../../Book1.tex
\section{一致可积性}
\label{b1:sec:0505}

% TODO: 撰写本节正文。

\subsection{一致可积族}
\label{b1:subsec:050501}

待填充。

\subsection{一致绝对连续性}
\label{b1:subsec:050502}

待填充。

\subsection{Vitali 收敛定理}
\label{b1:subsec:050503}

待填充。

\subsection{\texorpdfstring{\(L^1\)}{L1} 收敛判据}
\label{b1:subsec:050504}

待填充。


\begin{chaptersummary}
积分的构造保留了两层可加性：函数的非负线性组合可以拆开积分，积分区域的可数可测分割也可以拆开求和。从非负积分推广到一般函数时，正负部同时具有有限积分才给出可积性；这一条件让减法、复数线性和绝对误差估计都能合法进行。把几乎处处相等的代表合并后，积分绝对值成为 \(L^1\) 的范数。

\ProseParagraphBreak

收敛定理的适用范围由它们控制的对象决定。单调收敛沿着非负下逼近推进；Fatou 引理容许不受上界控制的非负函数列；控制收敛以一个可积函数同时限制全列，并给出 \(L^1\) 误差趋零。含参变量积分的求导还需要把导数控制转化为差商控制，不能只检查每个参数对应的积分存在。

\ProseParagraphBreak

本章按幅值尾部定义的一致可积性把共同控制函数放宽为统一的高值积分估计。在有限测度空间中，它与依测度收敛合起来恰好刻画 \(L^1\) 收敛；一般空间则还要控制有限测度集之外的积分。幅值尾部与空间尾部分别排除高度集中的尖峰和向更大区域分散的质量，不同文献是否把两者都纳入“一致可积”的定义，需要逐项核对。下一章将改变观察对象：固定一个函数对原测度加权会得到新测度，而绝对连续测度能否反过来由这样的函数表示，将由 Radon--Nikodym 定理回答。
\end{chaptersummary}

\BookExercises

\Cref{b1:ex:0509}与\cref{b1:ex:0510}集中训练三大收敛定理及三种收敛方式的基本辨析，适合在进入综合题以前完成；\cref{b1:ex:0503}至\cref{b1:ex:0508}需要组合本章若干工具，其中\cref{b1:ex:0508}可留待第二遍阅读。

% 由Axler §3A习题2的单个Dirac积分改编为可数原子、重复点与复值可积判据。
\begin{exercise}\label{b1:ex:0501}
设 \((X,\mathcal A)\) 是可测空间，\(x_k\in X\)、\(w_k\geq0\) 为有限实数，定义
\[
\mu(E)=\sum_{k=1}^{\infty}w_k\mathbf1_E(x_k).
\]
证明这是一个测度，且对任意非负可测函数 \(f\)，
\[
\int f\,\mathrm d\mu=\sum_{k=1}^{\infty}w_k f(x_k).
\]
允许点 \(x_k\) 重复，也不要求单点可测。给出复值函数关于 \(\mu\) 可积的相应判据。
\end{exercise}
\begin{solution}
对固定 \(x_k\)，映射 \(E\mapsto\mathbf1_E(x_k)\) 是测度，因为在一个互不相交的集合列中，\(x_k\) 至多属于其中一项。对两重非负级数交换求和次序，便知它们的加权和 \(\mu\) 可数可加。这里可通过对两个有限指标集分别取上确界来验证交换求和，不需要单点可测。

若 \(s=\sum_ja_j\mathbf1_{A_j}\) 非负且简单，则
\[
\int s\,\mathrm d\mu
=\sum_j a_j\sum_k w_k\mathbf1_{A_j}(x_k)
=\sum_k w_k s(x_k).
\]
取 \(s_n\uparrow f\)，左端由单调收敛趋于 \(\int f\,\mathrm d\mu\)；右端关于 \(n\) 递增，可以与非负求和交换极限，得到所求公式。将它用于 \(|f|\)，可知复值可测函数可积当且仅当
\(\sum_kw_k|f(x_k)|<\infty\)。此时对实部、虚部分别分解正负部，得到同样的积分求和公式；对应数值级数绝对收敛。
\end{solution}

% 自拟：把可测映射的逆像结构与本章积分构造相连接；不使用第7章换元公式。
\begin{exercise}\label{b1:ex:0502}
设 \(T:(X,\mathcal A,\mu)\rightarrow(Y,\mathcal B)\) 可测，令
\(\nu(B)=\mu\bigl(T^{-1}(B)\bigr)\)。证明 \(\nu\) 是测度，并对非负可测函数 \(h\) 证明
\[
\int_Yh\,\mathrm d\nu=\int_X h\circ T\,\mathrm d\mu.
\]
再说明同一公式如何推广到 \(\nu\)-可积的复值函数，及其 \(L^1\) 范数如何变化。
\end{exercise}
\begin{solution}
逆像保持空集、互不相交和可数并，所以 \(\mu\) 的测度公理直接给出 \(\nu\) 的测度公理。对 \(h=\mathbf1_B\)，所求公式就是 \(\nu\) 的定义；有限非负线性组合把公式推广到简单函数。一般非负 \(h\) 可取简单函数 \(s_n\uparrow h\)，于是 \(s_n\circ T\uparrow h\circ T\)，两端分别用单调收敛即可。

对可测复值函数应用刚证明的非负公式于 \(|h|\)，得到
\[
\|h\circ T\|_{L^1(\mu)}=\|h\|_{L^1(\nu)}.
\]
若右端有限，实部与虚部的正负分量都可积，把相应四个非负公式相减、相加即可。若 \(h_1=h_2\) 在一个 \(\nu\)-零集外成立，该零集的逆像为 \(\mu\)-零集，故复合也不依赖等价类代表。
\end{solution}

% 自拟广义控制收敛练习；变化的控制函数不能直接套用正文固定控制版本。
\begin{exercise}\label{b1:ex:0503}
设 \(f_n\to f\)、\(g_n\to g\) 几乎处处，其中 \(f_n,f\) 为实值或复值可测函数，\(g_n,g\geq0\) 可积。若
\[
|f_n|\leq g_n\quad\text{几乎处处},
\quad
\int g_n\,\mathrm d\mu\longrightarrow\int g\,\mathrm d\mu,
\]
证明 \(f\) 可积，且 \(\|f_n-f\|_1\to0\)。
\end{exercise}
\begin{solution}
在共同零集外取极限，有 \(|f|\leq g\)，故 \(f\) 可积。此时
\[
u_n=g_n+g-|f_n-f|\geq0
\]
几乎处处成立，且 \(u_n\to2g\)。在例外零集上统一改为零，再用 Fatou 引理，
\[
2\int g\,\mathrm d\mu
\leq\liminf_n\left(\int g_n\,\mathrm d\mu+\int g\,\mathrm d\mu-\|f_n-f\|_1\right)
=2\int g\,\mathrm d\mu-\limsup_n\|f_n-f\|_1.
\]
共同项有限，因而误差极限为零。证明没有断言 \((g_n)\) 被某一个可积函数逐点控制。
\end{solution}

% 自拟反例：用互不相交的小区间区分UI和共同可积控制。
\begin{exercise}\label{b1:ex:0504}
在 \((0,1)\) 上令
\[
E_n=\left(\frac1{n+1},\frac1n\right]\cap(0,1),
\quad f_n=n\mathbf1_{E_n}.
\]
证明 \(f_n\to0\) 处处且在 \(L^1\) 中成立，函数族一致可积，但不存在可积函数 \(g\) 使 \(|f_n|\leq g\) 几乎处处对所有 \(n\) 成立。
\end{exercise}
\begin{solution}
集合 \(E_n\) 互不相交，每点至多对应一个非零函数值，所以函数列处处趋零。并且
\[
\|f_n\|_1=n\left(\frac1n-\frac1{n+1}\right)=\frac1{n+1}\longrightarrow0.
\]
由\cref{b1:prop:l1-convergence-uniform-integrability}，该族一致可积；也可以直接估计其幅值尾部：
\(\sup_{n>M}\int f_n\leq1/(\lfloor M\rfloor+1)\to0\)。

若存在共同可积控制 \(g\)，先把所有控制不等式的例外零集取并。对于每个 \(r\)，不交性与单调性给出
\[
\int_0^1 |g|\,\mathrm dm\geq\sum_{n=1}^{r}\int_{E_n}f_n\,\mathrm dm
=\sum_{n=1}^{r}\frac1{n+1}.
\]
右端趋于无穷，与可积性矛盾。
\end{solution}

% 背景线索：Bass2024Real习题7.24（非正式引文）；增补无Hölder工具的定量小集估计。
\begin{exercise}\label{b1:ex:0505}
设 \(\alpha>0\)，可测函数族 \(\mathcal F\) 满足
\[
\sup_{f\in\mathcal F}\int_X|f|^{1+\alpha}\,\mathrm d\mu\leq C<\infty.
\]
对 \(0<\mu(A)<\infty\) 证明
\[
\sup_{f\in\mathcal F}\int_A|f|\,\mathrm d\mu
\leq(1+C)\mu(A)^{\alpha/(1+\alpha)}.
\]
再证明，当 \(\mu(X)<\infty\) 时，这族函数可积且一致可积。若 \(f_n\in\mathcal F\)，且 \(f_n\xrightarrow{\mu}f\)，其中 \(f\) 为可测实值或复值函数，证明 \(f\in L^1(\mu)\) 且 \(\|f_n-f\|_1\to0\)。不使用尚未建立的 Hölder 不等式。
\end{exercise}
\begin{solution}
给定 \(M>0\)，在 \(\{|f|\leq M\}\) 与其补集上分开估计：
\[
\int_A|f|\,\mathrm d\mu
\leq M\mu(A)+M^{-\alpha}\int_X|f|^{1+\alpha}\,\mathrm d\mu
\leq M\mu(A)+CM^{-\alpha}.
\]
令 \(M=\mu(A)^{-1/(1+\alpha)}\) 得到所需上界。若 \(\mu(A)=0\)，积分直接为零。有限测度时在 \(A=X\) 上得到可积性和一致的 \(L^1\) 上界；若全空间测度为零也直接成立。另一方面，
\[
\sup_{f\in\mathcal F}\int_{\{|f|>M\}}|f|\,\mathrm d\mu\leq CM^{-\alpha}\longrightarrow0,
\]
故一致可积。若 \(f_n\in\mathcal F\) 且 \(f_n\xrightarrow{\mu}f\)，则子族 \(\{f_n\}\) 也一致可积；由\cref{b1:thm:vitali-convergence}，\(f\in L^1(\mu)\) 且 \(\|f_n-f\|_1\to0\)。
\end{solution}

% 经典Frullani型积分的自编参数求导解法；仅使用本章已证Laplace零阶积分。
\begin{exercise}\label{b1:ex:0506}
设 \(a,b>0\)。证明函数
\[
x\longmapsto\frac{e^{-ax}-e^{-bx}}{x}
\]
在 \((0,\infty)\) 上可积，并计算其积分。要求说明零点附近、无穷远及参数求导各自所需的控制。
\end{exercise}
\begin{solution}
在 \(0<x\leq1\) 上，对参数使用中值定理，得到
\(|e^{-ax}-e^{-bx}|\leq|a-b|x\)。在 \(x\geq1\) 上，令 \(c=\min\{a,b\}\)，则绝对值除以 \(x\) 后不超过 \(2e^{-cx}\)。两段控制均可积，所以积分存在。

固定 \(b\)，记此积分为 \(F(a)\)。参数导数是 \(-e^{-ax}\)；在任何 \(a_0>0\) 的邻域 \(a_0/2<a<3a_0/2\) 中，导数由 \(e^{-a_0x/2}\) 控制。由\cref{b1:thm:differentiation-under-integral}及\cref{b1:prop:laplace-exponential-moments}，
\[
F'(a)=-\int_0^\infty e^{-ax}\,\mathrm dx=-\frac1a.
\]
因此 \(F(a)=-\log a+C_b\)。利用 \(F(b)=0\) 得 \(C_b=\log b\)，所求积分为 \(\log(b/a)\)。该计算不需要交换两个积分的顺序。
\end{solution}

% 自拟综合题：连接第3章正则性与本章简单函数L1逼近；不预借Tietze延拓或卷积。
\begin{exercise}\label{b1:ex:0507}
证明每个 \(f\in L^1(\mathbb R^d)\) 都能在 \(L^1\) 中由连续且在某个紧集外为零的函数逼近。可先处理有限测度可测集的示性函数：用紧集 \(K\) 和有界开集 \(O\supset K\) 逼近，再构造在 \(K\) 上为一、在 \(O\) 外为零的连续函数。
\end{exercise}
\begin{solution}
设 \(m(E)<\infty\)。由 Lebesgue 测度内正则性，给定 \(\varepsilon>0\)，可取紧集 \(K\subset E\)，使 \(m(E\setminus K)<\varepsilon\)。若 \(K\) 为空，零函数已给出所需误差。否则，由外正则性可取包含 \(K\) 的开集，再与包含 \(K\) 的大开球相交，得到有界开集 \(O\)，满足 \(m(O\setminus K)<\varepsilon\)。

以 \(\operatorname{dist}(x,A)=\inf_{y\in A}|x-y|\) 表示点到非空集合的距离，令
\[
\varphi(x)=
\frac{\operatorname{dist}(x,\mathbb R^d\setminus O)}
{\operatorname{dist}(x,\mathbb R^d\setminus O)+\operatorname{dist}(x,K)}.
\]
两个距离函数连续，且不会同时为零：若点到 \(K\) 的距离为零，则该点属于闭集 \(K\subset O\)，到闭集 \(\mathbb R^d\setminus O\) 的距离为正。因此 \(\varphi\) 连续，\(0\leq\varphi\leq1\)，在 \(K\) 上为一、在 \(O\) 外为零。它在紧集 \(\overline O\) 外为零，而且
\[
\|\mathbf1_E-\varphi\|_1\leq m(E\setminus K)+m(O\setminus K)<2\varepsilon.
\]

一般 \(f\) 先由可积简单函数 \(s=\sum_{j=1}^{r}a_j\mathbf1_{E_j}\) 逼近；舍去零系数后，每个 \(E_j\) 测度有限。对每个示性函数按上面方法选取 \(\varphi_j\)，使
\(\sum_j|a_j|\cdot\|\mathbf1_{E_j}-\varphi_j\|_1\) 任意小。有限和 \(\sum_ja_j\varphi_j\) 连续，且在有限个紧集的并之外为零。最后用三角不等式与 \(s\) 的逼近误差合并，便得到结论，复值系数也适用。
\end{solution}

% 自编证明：把UI的尾部预算组合成凸超线性函数，不引用高级Orlicz空间理论。
\begin{exercise}\label{b1:ex:0508}
设 \(\mathcal F\subset L^1(\mu)\) 非空。证明它按本章的幅值尾部定义一致可积，当且仅当存在有限值的非负递增凸函数
\(\Phi:[0,\infty)\rightarrow[0,\infty)\)，使
\[
\Phi(0)=0,\quad
\lim_{t\to\infty}\frac{\Phi(t)}t=\infty,\quad
\sup_{f\in\mathcal F}\int\Phi(|f|)\,\mathrm d\mu<\infty.
\]
构造方向可选 \(M_k\uparrow\infty\)，再考虑 \(\sum_k(t-M_k)_+\)。
\end{exercise}
\begin{solution}
若已有这样的 \(\Phi\)，记积分上界为 \(C\)，并令
\(c_M=\inf_{t>M}\Phi(t)/t\)。超线性增长保证 \(c_M\to\infty\)。当 \(M\) 足够大时，
\[
\sup_{f\in\mathcal F}\int_{\{|f|>M\}}|f|\,\mathrm d\mu
\leq C/c_M\longrightarrow0.
\]

反过来，利用一致可积性，递归选择严格递增且 \(M_k\geq k\) 的数，使
\[
\sup_{f\in\mathcal F}\int_{\{|f|>M_k\}}|f|\,\mathrm d\mu\leq2^{-k}.
\]
定义 \(\Phi(t)=\sum_{k=1}^{\infty}(t-M_k)_+\)。对每个有界的 \(t\) 区间只有有限项非零，所以函数有限、连续、递增且凸，并满足 \(\Phi(0)=0\)。固定正整数 \(r\)，当 \(t\geq2M_r\) 时前 \(r\) 项各不小于 \(t/2\)，于是 \(\Phi(t)/t\geq r/2\)。这证明超线性增长。

最后由非负级数积分公式，
\[
\int\Phi(|f|)\,\mathrm d\mu
=\sum_{k=1}^{\infty}\int(|f|-M_k)_+\,\mathrm d\mu
\leq\sum_{k=1}^{\infty}2^{-k}=1.
\]
同一个 \(\Phi\) 和同一个上界适用于全族，所以构造完成。该判据只控制幅值尾部，并不自动补上无穷测度空间中的集外控制条件。
\end{solution}

% 基础辨析自拟：要求逐项核对Fatou、MCT与DCT的假设，而不只计算积分。
\begin{exercise}\label{b1:ex:0509}
对下列函数列，分别判断 Fatou 引理、单调收敛定理和控制收敛定理的假设是否满足；若不满足，指出失败的假设。可以适用多个结果时，分别说明它们给出什么信息。
\begin{enumerate}
\item 在 \((0,1)\) 上取 \(f_n(x)=x^n\)；
\item 在 \(\mathbb R\) 上取 \(f_n=\mathbf1_{[0,n]}\)；
\item 在 \((0,1)\) 上取 \(f_n=n\mathbf1_{(0,1/n)}\)。
\end{enumerate}
\end{exercise}
\begin{solution}
第一列非负且处处趋于零，并由 \(\mathbf1_{(0,1)}\in L^1(0,1)\) 控制，所以控制收敛定理给出
\[
\int_0^1x^n\,\mathrm dx=\frac1{n+1}\longrightarrow0.
\]
Fatou 引理也适用，但这里只给出 \(0\leq0\)。该列递减而非递增，不能直接使用单调收敛定理。

第二列非负且递增至 \(\mathbf1_{[0,\infty)}\)，所以单调收敛定理给出
\[
\int_{\mathbb R}f_n\,\mathrm dx=n\uparrow\infty
=\int_{\mathbb R}\mathbf1_{[0,\infty)}\,\mathrm dx.
\]
Fatou 引理同样适用，但只给出相应的下半连续不等式。控制收敛定理不适用：若某个可积函数 \(g\) 几乎处处控制全部 \(f_n\)，则 \(g\geq1\) 几乎处处成立于 \([0,\infty)\)，不可能可积。

第三列非负且几乎处处趋于零，积分却恒为一。因此 Fatou 引理适用并给出严格不等式
\[
0=\int_0^1\liminf_n f_n\,\mathrm dx
<\liminf_n\int_0^1f_n\,\mathrm dx=1.
\]
这列既不递增，也不存在共同的可积控制函数；否则控制收敛定理会迫使积分趋于零。故后两个定理均不适用。这一项也说明，确认点态极限和逐项积分值并不能代替核对定理假设。
\end{solution}

% 基础收敛方式辨析自拟：同题比较依测度、几乎处处与L1收敛，并区分集中和逃逸。
\begin{exercise}\label{b1:ex:0510}
完成下列收敛方式比较。
\begin{enumerate}
\item 证明任意测度空间上的 \(L^1\) 收敛蕴含依测度收敛；再证明有限测度空间上的几乎处处收敛蕴含依测度收敛。
\item 在 \([0,1)\) 上，对 \(k\geq0\)、\(0\leq j<2^k\) 定义
\[
u_{2^k+j}=\mathbf1_{[j2^{-k},(j+1)2^{-k})}.
\]
证明 \(u_n\) 依测度收敛于零，但整列在每一点都不收敛。
\item 分别考察 \((0,1)\) 上的 \(h_n=n\mathbf1_{(0,1/n)}\) 与 \(\mathbb R\) 上的 \(v_n=\mathbf1_{[n,n+1]}\)。证明两列都几乎处处趋于零且 \(L^1\) 范数恒为一，判断它们是否依测度收敛，并说明哪一列体现质量集中，哪一列体现质量逃逸。
\end{enumerate}
\end{exercise}
\begin{solution}
若 \(\|f_n-f\|_1\to0\)，则对每个 \(\eta>0\)，由 Markov 型估计，
\[
\mu\{\,|f_n-f|>\eta\,\}
\leq\eta^{-1}\|f_n-f\|_1\longrightarrow0.
\]
若 \(\mu(X)<\infty\) 且 \(f_n\to f\) 几乎处处，则
\(\mathbf1_{\{|f_n-f|>\eta\}}\to0\) 几乎处处，并由可积函数 \(\mathbf1_X\) 控制。控制收敛定理给出这些示性函数的积分趋零，也就是依测度收敛。

对第二问，若 \(2^k\leq n<2^{k+1}\)，则 \(u_n\) 的支集测度为 \(2^{-k}\)，随 \(n\to\infty\) 趋于零。因此对 \(0<\eta<1\)，
\[
m\{\,|u_n|>\eta\,\}=2^{-k}\longrightarrow0,
\]
而 \(\eta\geq1\) 时偏差集为空，故 \(u_n\) 依测度趋零。另一方面，每个 \(x\in[0,1)\) 在每一级二进分割中恰好属于一个小区间，所以 \(u_n(x)=1\) 无穷多次；同一级其余区间又使 \(u_n(x)=0\)，故零也出现无穷多次。整列在任何一点都不收敛。

对第三问，每个固定点最终都离开相应支集，所以 \(h_n\to0\) 与 \(v_n\to0\) 都几乎处处成立，而且
\[
\|h_n\|_{L^1(0,1)}=1,
\qquad
\|v_n\|_{L^1(\mathbb R)}=1.
\]
对 \(0<\eta<1\)，充分大的 \(n\) 满足
\[
m\{\,|h_n|>\eta\,\}=\frac1n\longrightarrow0,
\qquad
m\{\,|v_n|>\eta\,\}=1.
\]
因此 \(h_n\) 依测度趋零而 \(v_n\) 不依测度趋零，两列都不在 \(L^1\) 中趋零。前者把固定质量压入越来越小的集合，体现质量集中；后者把固定质量移向无穷远，体现质量逃逸。
\end{solution}
